% !TeX root=../main.tex
\chapter{مروری بر مطالعات انجام شده}
%\thispagestyle{empty}

% ============================================================
%  مقدمه
% ============================================================
\section{مقدمه}
\label{subsec:lr_overview}

این فصل، نسب فکری پژوهش حاضر را در چهار حوزه‌ی دانشیِ مجاور ترسیم می‌کند:
مبانی مدیریت پروژه~(\lr{PM})، مدیریت و انتخاب سبد پروژه~(\lr{PPM})،
بودجه‌بندیِ سبد پروژه و مدل‌سازیِ جریان نقدی، و یادگیری تقویتی~(\lr{RL})
به‌کاررفته در مسائل تخصیصِ ترتیبیِ منابع. این چهار جریان در کنار هم، هم
بنیانِ نظریِ پژوهش و هم خلأیی را تعریف می‌کنند که انگیزه‌ی شکل‌گیریِ
چارچوب پیشنهادی است. این مرور از دستورالعمل‌های گزارش‌دهیِ
\lr{PRISMA~2020}~\citep{Page2021} پیروی می‌کند تا شفافیت و
بازتولیدپذیری در شناساییِ منابع، غربالگری، و تصمیماتِ گنجاندن را
تضمین کند.

% - - - - - - - - - - - - - - - - - - - - - - - - - - - - - -
\subsection{پروتکل جستجو و غربالگری}
\label{subsubsec:prisma}
% - - - - - - - - - - - - - - - - - - - - - - - - - - - - - -

منابع از پایگاه‌های \lr{Google Scholar}، \lr{Semantic Scholar} و
\lr{Scopus} با استفاده از راهبردی جستجو با محدودسازیِ تدریجی بازیابی
شدند. جستجوها از گسترده‌ترین حوزه‌ی مرتبط (\emph{مدیریت پروژه}) آغاز
شد و به‌ترتیب از طریق \emph{مدیریت سبد پروژه}،
\emph{بودجه‌بندیِ سبد و بهینه‌سازیِ جریان نقدی}، \emph{مدیریتِ ارزشِ
کسب‌شده تحت عدم‌قطعیت}، و \emph{یادگیریِ تقویتی برای زمان‌بندیِ پروژه و
تخصیصِ منابع} پالایش شد. پیمایشِ شبکه‌ی استنادها ---برف‌گلوله‌سازیِ
پس‌رو و پیش‌رو از مقالاتِ لنگرِ پراستناد--- برای بازیابیِ آثار بنیادینی
که با جستجوی کلیدواژه آشکار نشده بودند به‌کار رفت.

رکوردها در سه مرحله‌ی خودکارِ متوالی، منطبق با نمودار جریانِ
\lr{PRISMA~2020} (پیوست~\ref{app:prisma})، غربال شدند. خط‌لوله‌ی
غربالگری، به‌همراه پیکره‌ی بایگانی‌شده و خروجی‌های آن، به‌صورت یک
زنجیره‌ی پردازشیِ بازتولیدپذیر پیاده‌سازی شده است که محل مخزن و
نگاشتِ مفهومیِ آن در پیوست~\ref{app:repo_map} مستند شده است.

\medskip
\textbf{مرحله‌ی ۱ --- شناسایی و ماشین‌خوانایی.}\;
فایل‌های \lr{PDF} کاندیدا تحت تریاژِ کتاب‌شناختی و خوانایی به‌صورت
خودکار قرار گرفتند. هر سند از نظر بازده‌ی استخراجِ متن (حداقلِ تراکمِ
کاراکترِ خوانا در صفحات)، کیفیتِ واژگانی (نسبتِ توکن‌های الفبایی و
میانگین طول واژه‌ها)، و قابلیتِ جاسازی توسط مدلِ تبدیل‌گر جمله
ارزیابی شد. اسنادی که در هر یک از این معیارها ناموفق بودند در یک
بایگانیِ رد اختصاصی قرنطینه شدند، به‌گونه‌ای که هشِ محتواییِ هر سند
ثبت شد تا از ورود مجدد ارسال‌های تکراری به خط‌لوله جلوگیری شود.

\medskip
\textbf{مرحله‌ی ۲ --- فیلتر کردن بر اساس ارتباط.}\;
رکوردهای عبورکرده از مرحله‌ی ۱ بر پایه‌ی امتیازی مرکب از ارتباط
رتبه‌بندی شدند. یک مؤلفه‌ی شباهتِ معنایی، چکیده، مقدمه و نتیجه‌گیریِ
هر مقاله را با یک مدل تبدیل‌گر جمله جاسازی کرد و بیشینه‌ی شباهتِ
کسینوسی را نسبت به دوازده بردارِ پرسمانِ برگرفته از پایان‌نامه که
حوزه‌های دانشیِ هدف را پوشش می‌دادند محاسبه نمود. یک مؤلفه‌ی
امتیازدهیِ کلیدواژه‌ای نیز شمارش‌های وزن‌دارِ اصطلاحاتِ حوزه‌محور را
تخصیص داد. این دو مؤلفه با وزنِ یکسان ترکیب شدند تا امتیازِ ترکیبیِ
ارتباط در بازه‌ی $[0,1]$ به‌دست آید. مقالات به سه گروهِ \emph{بسیار
مرتبط} (امتیاز $\geq 0.45$)، \emph{ارزشمند} (امتیاز $\geq 0.30$)، و
\emph{نامرتبط} (امتیاز $< 0.30$) تریاژ شدند و دو گروهِ نخست برای
استخراجِ ساختاریافته نگه‌داشته شدند.

\medskip
\textbf{مرحله‌ی ۳ --- استخراج ساختاریافته.}\;
هر مقاله‌ی نگه‌داشته‌شده تحتِ استخراجِ دانشِ ساختاریافته قرار گرفت.
چکیده، مقدمه و نتیجه‌گیریِ هر \lr{PDF} تجزیه و به یک مدلِ زبانیِ
بزرگ با یک پرامپتِ طرحواره‌ایِ ثابت ارسال شد؛ خروجیِ این فرایند
رکوردی برای هر مقاله تولید می‌کرد که شاملِ موارد زیر بود: قاب‌بندیِ
مسئله، سطحِ سبد، نوعِ تصمیم، نحوه‌ی برخورد با جریانِ نقدیِ تصادفی،
سازوکارهای قراردادی، مقرراتِ تسویه‌ی فسخ، پارادایمِ مدل‌سازی،
پارامترهای کالیبراسیون، و پرچم‌های ارتباط با پایان‌نامه. پیکره‌ی
ساختاریافته‌ی حاصل، یک داشبوردِ سنتزِ خودکار و یک مولدِ جدول را هدایت
کرد که ماتریسِ تحلیلِ شکاف، نمودارهای توزیعِ حوزه‌ها، و نقشه‌ی ادبیاتِ
بازتولیدشده در جدول~\ref{tab:litmap} را تولید نمود.

\medskip
خط‌لوله‌ی کامل ---اسکریپت‌های پردازشی، پیکره‌ی رکوردهای ساختاریافته،
خروجی‌های داشبورد، و جداول سنتز--- در مخزنِ پروژه بایگانی شده و یک
پیکره‌ی مرورِ ادبیاتِ بازتولیدپذیر و قابل‌حسابرسیِ ماشینی را تشکیل
می‌دهد؛ سازمان‌دهیِ سطحِ‌دایرکتوریِ آن در پیوست~\ref{app:repo_map}
شرح داده شده است. نمودار جریانِ \lr{PRISMA} که رکوردهای هر مرحله را
کمّی می‌کند در پیوست~\ref{app:prisma} ارائه شده است.

\begin{figure}[htbp]
    \centering
    % figures/screening_pipeline.tex
% TikZ only — no figure environment, no caption, no label.
% Include in section file as:
%   % figures/screening_pipeline.tex
% TikZ only — no figure environment, no caption, no label.
% Include in section file as:
%   % figures/screening_pipeline.tex
% TikZ only — no figure environment, no caption, no label.
% Include in section file as:
%   \input{figures/screening_pipeline.tex}
\begin{tikzpicture}[
    box/.style      = {rectangle, rounded corners=6pt, thick,
                       text width=7.0cm, align=center,
                       minimum height=0.85cm, inner sep=7pt,
                       font=\small},
    decision/.style = {diamond, thick, aspect=2.5,
                       minimum width=4.5cm, minimum height=0.7cm,
                       align=center, font=\scriptsize, inner sep=2pt},
    rej/.style      = {rectangle, rounded corners=6pt, thick,
                       text width=3.2cm, align=center,
                       minimum height=0.75cm, inner sep=5pt,
                       font=\scriptsize},
  ]

  %% --- Row 1 : left -> right ---
  \node[box, flowColorA] (s1)
    {\textbf{Stage 1 --- Identification \& Readability}\\
     Extraction yield $\cdot$ Lexical quality $\cdot$ Embeddability};

  \node[decision, flowColorGray, right=1.2cm of s1] (d1)
    {All three\\tests pass?};

  %% --- Row 2 : right -> left ---
  \node[box, flowColorB, below=2.4cm of d1] (s2)
    {\textbf{Stage 2 --- Relevance Filtering}\\
     Semantic similarity $+$ Keyword score $\to r \in [0,1]$};

  \node[decision, flowColorGray, left=1.2cm of s2] (d2)
    {$r \geq 0.30$?};

  %% --- Row 3 : left -> right ---
  \node[box, flowColorC, below=1.7cm of d2] (s3)
    {\textbf{Stage 3 --- Structured Extraction}\\
     JSON extraction: metadata $\cdot$ calibration $\cdot$ thesis flags};

  \node[box, flowColorD, right=1.2cm of s3] (out)
    {\textbf{Synthesis}\\
     Gap matrix $\cdot$ Domain charts $\cdot$ Literature map};

  %% --- Rejection nodes ---
  \node[rej, flowColorReject, above=0.55cm of d1] (r1)
    {Rejected\\(failed readability)};

  \node[rej, flowColorReject, above=0.55cm of d2] (r2)
    {Rejected\\(irrelevant)};

  %% --- Main snake flow ---
  \draw[arrowMain] (s1) -- (d1);
  \draw[arrowMain] (d1) -- node[right, font=\scriptsize]{Yes} (s2);
  \draw[arrowMain] (s2) -- (d2);
  \draw[arrowMain] (d2) -- node[left,  font=\scriptsize]{Yes} (s3);
  \draw[arrowMain] (s3) -- (out);

  %% --- Rejection branches ---
  \draw[arrowReject] (d1) -- node[right, font=\scriptsize]{No} (r1);
  \draw[arrowReject] (d2) -- node[left,  font=\scriptsize]{No} (r2);

\end{tikzpicture}
\begin{tikzpicture}[
    box/.style      = {rectangle, rounded corners=6pt, thick,
                       text width=7.0cm, align=center,
                       minimum height=0.85cm, inner sep=7pt,
                       font=\small},
    decision/.style = {diamond, thick, aspect=2.5,
                       minimum width=4.5cm, minimum height=0.7cm,
                       align=center, font=\scriptsize, inner sep=2pt},
    rej/.style      = {rectangle, rounded corners=6pt, thick,
                       text width=3.2cm, align=center,
                       minimum height=0.75cm, inner sep=5pt,
                       font=\scriptsize},
  ]

  %% --- Row 1 : left -> right ---
  \node[box, flowColorA] (s1)
    {\textbf{Stage 1 --- Identification \& Readability}\\
     Extraction yield $\cdot$ Lexical quality $\cdot$ Embeddability};

  \node[decision, flowColorGray, right=1.2cm of s1] (d1)
    {All three\\tests pass?};

  %% --- Row 2 : right -> left ---
  \node[box, flowColorB, below=2.4cm of d1] (s2)
    {\textbf{Stage 2 --- Relevance Filtering}\\
     Semantic similarity $+$ Keyword score $\to r \in [0,1]$};

  \node[decision, flowColorGray, left=1.2cm of s2] (d2)
    {$r \geq 0.30$?};

  %% --- Row 3 : left -> right ---
  \node[box, flowColorC, below=1.7cm of d2] (s3)
    {\textbf{Stage 3 --- Structured Extraction}\\
     JSON extraction: metadata $\cdot$ calibration $\cdot$ thesis flags};

  \node[box, flowColorD, right=1.2cm of s3] (out)
    {\textbf{Synthesis}\\
     Gap matrix $\cdot$ Domain charts $\cdot$ Literature map};

  %% --- Rejection nodes ---
  \node[rej, flowColorReject, above=0.55cm of d1] (r1)
    {Rejected\\(failed readability)};

  \node[rej, flowColorReject, above=0.55cm of d2] (r2)
    {Rejected\\(irrelevant)};

  %% --- Main snake flow ---
  \draw[arrowMain] (s1) -- (d1);
  \draw[arrowMain] (d1) -- node[right, font=\scriptsize]{Yes} (s2);
  \draw[arrowMain] (s2) -- (d2);
  \draw[arrowMain] (d2) -- node[left,  font=\scriptsize]{Yes} (s3);
  \draw[arrowMain] (s3) -- (out);

  %% --- Rejection branches ---
  \draw[arrowReject] (d1) -- node[right, font=\scriptsize]{No} (r1);
  \draw[arrowReject] (d2) -- node[left,  font=\scriptsize]{No} (r2);

\end{tikzpicture}
\begin{tikzpicture}[
    box/.style      = {rectangle, rounded corners=6pt, thick,
                       text width=7.0cm, align=center,
                       minimum height=0.85cm, inner sep=7pt,
                       font=\small},
    decision/.style = {diamond, thick, aspect=2.5,
                       minimum width=4.5cm, minimum height=0.7cm,
                       align=center, font=\scriptsize, inner sep=2pt},
    rej/.style      = {rectangle, rounded corners=6pt, thick,
                       text width=3.2cm, align=center,
                       minimum height=0.75cm, inner sep=5pt,
                       font=\scriptsize},
  ]

  %% --- Row 1 : left -> right ---
  \node[box, flowColorA] (s1)
    {\textbf{Stage 1 --- Identification \& Readability}\\
     Extraction yield $\cdot$ Lexical quality $\cdot$ Embeddability};

  \node[decision, flowColorGray, right=1.2cm of s1] (d1)
    {All three\\tests pass?};

  %% --- Row 2 : right -> left ---
  \node[box, flowColorB, below=2.4cm of d1] (s2)
    {\textbf{Stage 2 --- Relevance Filtering}\\
     Semantic similarity $+$ Keyword score $\to r \in [0,1]$};

  \node[decision, flowColorGray, left=1.2cm of s2] (d2)
    {$r \geq 0.30$?};

  %% --- Row 3 : left -> right ---
  \node[box, flowColorC, below=1.7cm of d2] (s3)
    {\textbf{Stage 3 --- Structured Extraction}\\
     JSON extraction: metadata $\cdot$ calibration $\cdot$ thesis flags};

  \node[box, flowColorD, right=1.2cm of s3] (out)
    {\textbf{Synthesis}\\
     Gap matrix $\cdot$ Domain charts $\cdot$ Literature map};

  %% --- Rejection nodes ---
  \node[rej, flowColorReject, above=0.55cm of d1] (r1)
    {Rejected\\(failed readability)};

  \node[rej, flowColorReject, above=0.55cm of d2] (r2)
    {Rejected\\(irrelevant)};

  %% --- Main snake flow ---
  \draw[arrowMain] (s1) -- (d1);
  \draw[arrowMain] (d1) -- node[right, font=\scriptsize]{Yes} (s2);
  \draw[arrowMain] (s2) -- (d2);
  \draw[arrowMain] (d2) -- node[left,  font=\scriptsize]{Yes} (s3);
  \draw[arrowMain] (s3) -- (out);

  %% --- Rejection branches ---
  \draw[arrowReject] (d1) -- node[right, font=\scriptsize]{No} (r1);
  \draw[arrowReject] (d2) -- node[left,  font=\scriptsize]{No} (r2);

\end{tikzpicture}
    \caption{خط‌لوله‌ی غربالگریِ سه‌مرحله‌ایِ منطبق با \lr{PRISMA}، به‌صورت
      جریانی مارپیچی چیده شده است. جعبه‌ها نشان‌دهنده‌ی مراحل پردازشِ
      خودکار؛ لوزی‌ها نشان‌دهنده‌ی دروازه‌های قبول/رد؛ و فلش‌های
      نقطه‌چین نشان‌دهنده‌ی مسیرِ رد کردن هستند. جزئیاتِ پیاده‌سازی و
      محلِ مخزن در پیوست~\ref{app:repo_map} ارائه شده است.}
    \label{fig:screening_pipeline}
\end{figure}

% ============================================================
%  تعاریف، اصول و مبانی نظری
% ============================================================
\section{تعاریف، اصول و مبانی نظری}
\label{subsec:lr_pm}

مطالعه‌ی آکادمیکِ مدیریت پروژه بر سه ستونِ کمّیِ اصلیِ کنترلِ
زمان‌بندی (\lr{CPM}/\lr{PERT})، مدیریتِ ارزشِ کسب‌شده~(\lr{EVM})، و
مدل‌سازیِ ریسک استوار است~\citep{Fleming2010, PMI2021}. \lr{EVM}
شاخصِ عملکردِ هزینه و شاخصِ عملکردِ زمان‌بندی را به‌عنوانِ سیگنال‌های
تشخیصیِ بلادرنگ معرفی کرد که هزینه‌ی بودجه‌ایِ کارِ انجام‌شده را با
مخارجِ واقعی و زمان‌بندیِ برنامه‌ریزی‌شده مرتبط می‌سازند؛ این دو شاخص
واژگانِ عملکردیِ پایه‌ای را فراهم می‌کنند که مدل حاضر مستقیماً از
آن‌ها بهره می‌برد. مدل حاضر صورت‌بندیِ «زمان‌بندیِ کسب‌شده» از شاخصِ
عملکردِ زمان‌بندی را به‌کار می‌گیرد~\citep{Lipke2003}؛ این صورت‌بندی،
برخلافِ نسبتِ کلاسیک، پس از تاریخِ پایانِ برنامه‌ریزی‌شده نیز همچنان
کاهش می‌یابد و به‌محضِ رسیدنِ ارزشِ برنامه‌ریزی‌شده به بودجه‌ی کل در
لحظه‌ی تکمیل منجمد نمی‌شود ---محدودیتی شناخته‌شده در نسبتِ کلاسیک طیِ
دوره‌های بازیابی پس از پایانِ برنامه‌ریزی‌شده~\citep{Lipke2003,
Henderson2004}. شاخصِ عملکرد تا تکمیل نیز به‌عنوانِ مکملی آینده‌نگر
در نظر گرفته شده و کاراییِ هزینه‌ایِ لازم در کار باقی‌مانده برای
پایان‌یافتنِ پروژه در چارچوبِ بودجه را اندازه می‌گیرد. برآورد در
تکمیل از مخرجی مرکب از هزینه و زمان‌بندی استفاده می‌کند؛
\citet{christensen1993cost} با بررسیِ ۱۵۵ قرارداد نشان داد که این
صورت‌بندی دقیق‌ترین و پایدارترین پیش‌بینی‌های هزینه‌ی در تکمیل را
ارائه می‌دهد و اتکاپذیریِ آن از نقطه‌ی ۲۰ درصدِ تکمیل به بعد افزایش
می‌یابد~\citep{Lipke2003}. توسعه‌های تصادفیِ \lr{EVM}~\citep{Acebes2015,
Batselier2017} و چارچوب‌های شبیه‌سازیِ مونت‌کارلو~\citep{Kwak2007}
عدم‌قطعیتِ عملکرد را در سطحِ پروژه مدل‌سازی کرده و نشان داده‌اند که
شاخص‌های عملکرد از توزیع‌های تجربیِ قابلِ‌مشخصه‌سازی در طولِ عمرِ پروژه
پیروی می‌کنند؛ این یافته‌ها مبنای کالیبراسیونِ پارامترِ بازدهیِ تصادفی
در محیط حاضر را فراهم می‌کنند.

تحلیلِ منحنیِ \lr{S}، که توسط \citet{Kenley2009} برای جریان‌های نقدیِ
ساخت‌وساز رسمی‌سازی شد، نشان داد که مصرفِ بودجه از پروفیل‌هایِ
سیگموئیدیِ مشخصی پیروی می‌کند که پارامترهای شکلِ آن به نوعِ پروژه و
فرمِ قرارداد وابسته است. تابعِ بتای ناقصِ نرمال‌شده، پارامتری‌سازیِ
انعطاف‌پذیری تک‌خانواده‌ای از این پروفیل فراهم
می‌کند~\citep{barraza2004probabilistic} و مستقیماً در مدلِ مخارجِ
برنامه‌ریزی‌شده‌ی صورت‌بندیِ حاضر اتخاذ شده است.

شواهدِ تجربیِ بزرگ‌مقیاس درباره‌ی عملکردِ پروژه توسط \citet{Merrow2011}
از طریقِ پایگاه‌داده‌ی تحلیلِ مستقلِ پروژه، و توسط
\citet{Flyvbjerg2003, Flyvbjerg2018} از طریقِ مجموعه‌داده‌های
افزایشِ هزینه در پروژه‌های زیرساختی در دهه‌ها و جغرافیاهای مختلف
فراهم شده است. این مجموعه‌داده‌ها هم کالیبراسیونِ توزیع‌های عدم‌قطعیت
و هم معیارهایی که پارامترهایِ محیطِ مصنوعی در برابرِ آن‌ها
اعتبارسنجی می‌شوند را شکل می‌دهند. معیارهای تجربیِ تکمیلی نیز از
پایگاه‌دادهٔ‌های مهندسیِ هزینه‌ی \lr{AACE International} و
\lr{Construction Industry Institute} استخراج شده‌اند~\citep{christensen1993cost}.

ساختارهای حاکمیتِ قراردادی ---به‌ویژه فرم‌های \lr{FIDIC} (کتاب‌های
قرمز، زرد و نقره‌ای)--- تعاملِ مالیِ میانِ کارفرما و پیمانکار را از
طریقِ برنامه‌های پرداختِ نقطه‌ی عطف، سازوکارهای پیش‌پرداخت و بازیابیِ
آن، و بندهای حسنِ‌انجامِ‌کار تجویز می‌کنند~\citep{FIDIC2017}. طبقِ
بندِ فرعیِ ۱۴.۲ \lr{FIDIC}، پیش‌پرداخت وامی بدونِ بهره برای بسیجِ
کارگاه است که از طریقِ کسورِ درصدی از گواهی‌های پرداختِ بین‌دوره‌ای
بازیابی می‌شود؛ آستانه‌ی شروع و نرخِ بازیابی به‌عنوانِ شرایطِ خصوصیِ
پروژه‌محور مشخص می‌شوند~\citep{FIDIC2017}. مدل حاضر این سازوکار را با
نرخ‌های پیش‌پرداخت و بازیابیِ خاصِ هر پروژه پارامتری‌سازی می‌کند، و بدین
ترتیب طیفِ کاملِ شرایطِ خصوصیِ رایج در عمل را در برمی‌گیرد.

% ============================================================
%  مروری بر ادبیات موضوع
% ============================================================
\section{مروری بر ادبیات موضوع}

% ------------------------------------------------------------
\subsection{مدیریت سبد پروژه: انتخاب و حاکمیت}
\label{subsec:lr_ppm}
% ------------------------------------------------------------

مدیریتِ سبدِ پروژه، مسئله‌ی تخصیصِ منابعِ محدودِ یک سازمان را در میانِ
مجموعه‌ای رقیب از پروژه‌های کاندیدا برای بیشینه‌سازیِ ارزشِ راهبردی
رسمی‌سازی می‌کند~\citep{Archer1999, Cooper2001}. ادبیاتِ موضوع دو
لایه‌ی تصمیم‌گیریِ تودرتو را از هم متمایز می‌کند: \emph{انتخابِ سبد}
---بهینه‌سازیِ دودویی یا عددِ‌صحیحِ‌مختلط برایِ تعیینِ پروژه‌هایی که
باید گنجانده شوند--- و \emph{حاکمیتِ سبد}---نظارتِ عملیاتیِ مستمر بر
پروژه‌های تأییدشده و فعال.

مدل‌های انتخابِ سبد به‌طور گسترده توسعه یافته‌اند. چارچوب‌های
میانگین‌واریانسِ سبکِ مارکویتز که از نظریه‌ی سبدِ مالی اقتباس شده‌اند
\citep{Markowitz1952}، برای برقراریِ موازنه میانِ ارزشِ فعلیِ خالصِ
موردِ انتظار و واریانس، بر سبدهای پروژه اعمال شده‌اند~\citep{Eilat2006}.
توسعه‌های تحلیلِ تصمیم‌گیریِ چندمعیاره، همسوییِ راهبردی، رتبه‌بندیِ
ریسک، و وابستگی‌های میانِ منابع را نیز در برمی‌گیرند~\citep{Martinsuo2007,
Killen2008}. رویکردهای فراابتکاری برای نمونه‌های ترکیباتیِ بزرگ
به‌کار رفته‌اند~\citep{Zhang2020}. انتخابِ سبدِ مقیدشده به منابع تحتِ
سهمیه‌بندیِ سرمایه توسط \citet{kavadias2003resource} مطالعه شده است؛
قیدِ عدمِ‌استقراضِ او مستقیماً در صورت‌بندیِ حاضر اتخاذ شده است.

پژوهشِ حاضر، خروجیِ انتخابِ سبد را صراحتاً به‌عنوانِ ورودیِ برون‌زا
می‌پذیرد. مسئله‌ی انتخاب، حل‌شده در بالادست فرض می‌شود؛ تمرکزِ کار
به‌طورِ کامل بر لایه‌ی حاکمیتِ عملیاتی است که سبدی متعهدشده را
دریافت می‌کند و باید آن را در طولِ اجرا مدیریت کند. این تجزیه با
مدلِ حاکمیتِ دومرحله‌ایِ \citet{Turner2014} و استانداردِ \lr{PMI} برای
مدیریتِ سبد~\citep{PMI2017} سازگار است که انتخابِ راهبردی را از
نظارتِ اجراییِ عملیاتی جدا می‌کند. بدین‌ترتیب، این پژوهش مستقیماً بر
ادبیاتِ انتخاب بنا می‌شود و در عینِ‌حال مرزِ دانش را به حوزه‌ی
حاکمیت می‌راند که توجهِ تحلیلیِ رسمیِ به‌مراتب کمتری دریافت کرده است.

% ------------------------------------------------------------
\subsection{بودجه‌بندیِ سبد پروژه، مدل‌سازیِ جریان نقدی و عدم‌قطعیت}
\label{subsec:lr_cashflow}
% ------------------------------------------------------------

در درونِ حاکمیتِ سبد، تخصیصِ بودجه میانِ پروژه‌های فعال در طولِ زمان
یک تصمیمِ عملیاتیِ پرتکرار و پرمخاطره را نمایندگی می‌کند. با وجودِ
اهمیتِ عملیِ آن، این مسئله‌ی خاص نسبت به انتخابِ سبد توجهِ رسمیِ
محدودتری دریافت کرده است. آثارِ موجود، مدل‌سازیِ جریانِ نقدی در سطحِ
پروژه~\citep{Kenley2009, Park2004}، پیش‌بینیِ هزینه‌ی در تکمیلِ
پروژه‌ی منفرد بر مبنایِ \lr{EVM}~\citep{Kim2002, Acebes2015}، و
تراز‌سازیِ منابع در پروژه‌های چندگانه~\citep{Lova2000, Peteghem2010}
را پوشش می‌دهند. این آثار بودجه‌ها را قیودی قطعی تلقی می‌کنند، نه
متغیرهای تصمیمِ پویا که تابعِ بهینه‌سازیِ ترتیبی در سطحِ سبد باشند.

عدم‌قطعیتِ جریانِ نقدی در مدلِ حاضر از یک منبعِ واحد ناشی می‌شود:
بازدهیِ بودجه‌به‌پیشرفت، که تعیین می‌کند یک تخصیصِ بودجه‌ی معین در هر
دوره چه میزان پیشرفتِ فیزیکی ایجاد
می‌کند~\citep{Acebes2015, Batselier2017}. بهره‌وریِ نامطمئن تعیین
می‌کند که آیا و چه زمانی هر یک از آستانه‌های نقطه‌ی عطف رد می‌شود، و
این امر به‌نوبه‌ی خود زمان‌بندیِ پرداختِ ثابتِ قراردادیِ نقطه‌ی عطف را
تعیین می‌کند. تأخیرهای پرداخت میانِ تأیید و دریافت به‌عنوانِ منبعِ
مستقلی از عدم‌قطعیت مدل نمی‌شوند؛ کانالِ عملکردمحور، منبعِ ساختاریِ
اصلیِ عدم‌قطعیتِ زمانیِ جریانِ نقدی در قراردادهایِ پرداختِ نقطه‌ی عطف
است، و جداسازیِ آن شناساییِ شفاف‌تری از مسئله‌ی تخصیص را بدونِ از
دست‌دادنِ قاب‌بندیِ اصلیِ عدم‌قطعیت ممکن می‌سازد. سایرِ منابعِ
عدم‌قطعیتِ جریانِ نقدی ---تغییراتِ محدوده‌ی کار، شوک‌های بازارِ
خارجی~\citep{Doloi2012}--- برون‌زا در نظر گرفته می‌شوند.

برنامه‌ریزیِ تصادفی~\citep{Birge2011} چارچوبِ کانونیِ تحقیق در
عملیات برای بهینه‌سازی تحتِ چنین عدم‌قطعیتی را فراهم می‌کند، و نوعِ
دومرحله‌ایِ آن برای زمان‌بندیِ پروژه تحتِ عدم‌قطعیتِ منابع به‌کار رفته
است~\citep{Lova2000}. با این‌ حال، برنامه‌های تصادفیِ چندمرحله‌ای برای
تخصیصِ بودجه‌ی ترتیبی در سطحِ سبد، با افزایشِ تعدادِ پروژه‌ها و
دوره‌ها با انفجارِ درختِ سناریو مواجه می‌شوند، همان‌گونه که در
بخش~\ref{subsec:theoretical_ub} رسمی‌سازی شده است. همین دشواری،
صورت‌بندیِ \lr{MDP} در پژوهشِ حاضر را توجیه می‌کند.

خلأِ موجود در ادبیات دقیق و مشخص است: هیچ چارچوبِ تثبیت‌شده‌ای
تصمیمِ تخصیصِ بودجه در سطحِ سبد را به‌عنوانِ مسئله‌ای از بهینه‌سازیِ
تصادفیِ ترتیبی مدل نمی‌کند که در آن یک تصمیم‌گیرنده‌ی مرکزی منابع را
میانِ چند پروژه‌ی فعال تخصیص می‌دهد، پروژه‌هایی که مسیرِ ورودیِ
نقدیِ آن‌ها به‌واسطه‌ی بازدهیِ تصادفیِ بودجه‌به‌پیشرفت نامطمئن است،
آن‌هم تحتِ پنجره‌های قراردادیِ تأییدِ نقطه‌ی عطف که سرعتِ تبدیلِ
پیشرفت به ورودیِ نقدی را محدود می‌کنند.

تابعِ پاداشِ محیطِ حاضر، جریانِ نقدیِ خالصِ تنزیل‌شده‌ی خالص است؛ هیچ
جریمه‌ی انحرافی در هدف وارد نمی‌شود. این انتخابی طراحی‌شده و عمدی
است: پیامدهایِ انحرافِ پایدار ---فسخ و از دست‌رفتنِ تمامیِ پرداخت‌های
آتیِ نقطه‌ی عطف و مبالغِ نگه‌داشته‌شده--- برای نامطلوب‌کردنِ انحراف
بدونِ نیاز به جمله‌ی جریمه‌ای جداگانه کافی‌اند، و گنجاندنِ چنین جمله‌ای
وزنِ اسکالرسازیِ دلبخواهانه‌ای را معرفی می‌کرد که هیچ کالیبراسیونِ
طبیعی‌ای از داده‌های قراردادی ندارد~\citep{Ng1999}. ضریبِ تنزیل، نه
فقط ترجیحِ زمانی، بلکه هزینه‌ی فرصتِ سرمایه‌ی درگیرشده در محیطی
رقابتی را نیز رمزگذاری می‌کند که در آن، پروژه‌های آینده در حالی که
پروژه‌های جاری ناتمام مانده‌اند از دست می‌روند~\citep{Damodaran2012}؛
همین امر لنگرِ کالیبراسیونِ مشخصی از نرخ‌های استقرارِ مجددِ خطِ‌لوله‌ی
مشاهده‌شده‌ی شرکت یا معیارهای رایجِ هزینه‌ی سرمایه به آن می‌دهد.
پیامدِ این طراحی آن است که یک سیاستِ خوش‌عملکرد ممکن است مسیرهایِ
تخصیصی تولید کند که مخارج را حولِ آستانه‌های تأییدِ نقطه‌ی عطف
متمرکز می‌سازد، و این دقیقاً بازتابِ راهبردِ مالیِ مستندِ پیمانکاران
در قراردادهایِ پرداختِ نقطه‌ی عطف است~\citep{Kenley2009, Kaka1993,
Arditi2005}.

% ------------------------------------------------------------
\subsection{عملکرد ضعیف پیمانکار، دوره‌های مهلت و فسخ قرارداد}
\label{subsec:lr_interventions}
% ------------------------------------------------------------

هنگامی که پیمانکاری در رعایتِ آستانه‌های عملکردِ زمان‌بندی یا هزینه
ناکام می‌ماند، قراردادهای بین‌المللیِ \lr{EPC} توالیِ تشدیدِ
ساختاریافته‌ای را پیش از امکانِ فسخ تجویز می‌کنند. طبقِ بندِ فرعیِ
۱۵.۱ \lr{FIDIC}، کارفرما باید اخطارِ رسمیِ نقض صادر کند و دوره‌ای
مهلت ---معمولاً ۲۸ روز--- اعطا کند که طیِ آن پیمانکار موظف است
بهبودِ قابلِ‌اندازه‌گیری را نشان دهد پیش از آنکه کارفرما بتواند به
فسخ اقدام کند~\citep{FIDIC2017, Ndekugri2008}. این دوره‌ی مهلت امتیازی
اختیاری نیست؛ تعهدی قراردادی است که نادیده‌گرفتنِ آن فسخِ ناروا
محسوب می‌شود. مطالعاتِ تجربیِ اختلافاتِ بین‌المللیِ \lr{EPC} تأیید
می‌کنند که عدمِ رعایتِ دوره‌ی مهلت یکی از پرتکرارترین مبانیِ موردِ
نزاع در ادعاهای فسخِ پیمانکار است~\citep{Ndekugri2008, Ennis2021}.

ادبیاتِ مدیریتِ پروژه‌ی رفتاری، شرایطی را مطالعه کرده که در آن‌ها
سازمان‌ها در خاتمه‌دادن به پروژه‌های کم‌عملکرد در نقطه‌ی اقتصادیِ
عقلانی ناکام می‌مانند. تشدیدِ تعهد~\citep{Staw1987}---گرایش به ادامه‌ی
سرمایه‌گذاری در پروژه‌های روبه‌شکست، فراتر از آنچه ارزشِ موردِانتظار
توجیه می‌کند--- در برنامه‌های ساخت‌وساز، فناوریِ اطلاعات و زیرساخت
مستند شده است. نظریه‌ی اختیاراتِ واقعی~\citep{Dixit1994,
Huchzermeier2001} مکملِ هنجاریِ آن را فراهم می‌کند: فسخ، اعمالِ یک
اختیارِ رهاسازی است که ارزشِ آن، پس از عبورِ عملکرد از آستانه‌ای
قابلِ‌بازیابی، از ارزشِ فعلیِ خالصِ ادامه فراتر می‌رود. این دو جریان
در کنارِ هم نشان می‌دهند که تصمیم به ادامه یا فسخِ پروژه‌ای کم‌عملکرد،
نه صرفاً قراردادی است و نه صرفاً مالی، بلکه هم‌زمان سیگنال‌های
عملکرد، تعهداتِ قراردادی، و زمان‌بندی را درگیر می‌کند.

مدل حاضر هر سه بُعد را در یک سازوکارِ محیطیِ واحد، تحتِ حاکمیتِ دو
شرط، ادغام می‌کند. نقضِ تأخیرِ پیشرفت زمانی رخ می‌دهد که فاصله‌ی
میانِ پیشرفتِ برنامه‌ریزی‌شده و واقعی از تحملِ قراردادی فراتر رود؛
این شرط بر پایه‌ی اختیارِ مهندسِ \lr{FIDIC} برای صدورِ اخطار هنگامی
است که پیشرفت به‌طورِ محسوسی از برنامه عقب می‌افتد~\citep{FIDIC2017}.
نقضِ مهلتِ سخت نیز بدونِ هیچ شرطی در تجاوز از پایانِ مجازِ قراردادی
رخ می‌دهد و فسخِ فوری را صرفِ‌نظر از وضعیتِ شمارنده‌ی دوره‌ی مهلت
فعال می‌کند. هنگامی که نقضِ تأخیرِ پیشرفت فعال باشد، شمارنده‌ی
دوره‌ی مهلت یک واحد کاهش می‌یابد؛ این شمارنده تنها زمانی که هیچ نقضی
وجود ندارد به مقدارِ اولیه‌ی خود بازمی‌گردد. فسخ هنگامی اجرا می‌شود
که شمارنده به صفر برسد. این سازوکار هم‌زمان در حقوقِ قراردادیِ
\lr{FIDIC}~\citep{FIDIC2017, Ndekugri2008}، قاب‌بندیِ اختیاراتِ واقعی
از آستانه‌ی رهاسازی~\citep{Dixit1994, Huchzermeier2001}، اعتبارسنجیِ
ادبیاتِ \lr{EVM} از شاخص‌های عملکرد به‌عنوانِ سیگنال‌هایِ پایدار و
پیشرو~\citep{christensen1993est, Lipke2003}، و شواهدِ ادبیاتِ تشدیدِ
تعهد مبنی بر اینکه سیگنال‌های عملکرد باید پایدار ---و نه گذرا--- باشند
پیش از آنکه فسخ موجه باشد~\citep{Staw1987}، ریشه دارد.

شمارنده‌ی دوره‌ی مهلت بخشی از حالتِ قابلِ‌مشاهده است و به سیاست
سیگنالی پیش‌دستانه می‌دهد که واقعیتِ قراردادی را بازتاب می‌دهد، بدونِ
آنکه کنترلِ اختیاری بر نتیجه‌ی فسخ به آن اعطا شود. از منظرِ مدیرِ
سبد، این سیگنال کارکردی دوگانه دارد: از یک سو پروژه‌هایی را که در
آستانه‌ی شکستِ قراردادی قرار دارند علامت می‌زند، و از سویِ دیگر نقطه‌ای
را مشخص می‌کند که از آن پس، تخصیصِ مستمر بازدهیِ کاهشی‌ای نسبت به
ارزشِ تسویه‌ی کارِ از پیش انجام‌شده دارد. مدیرِ عقلانی‌ای که شمارنده‌ای
را در حالِ نزدیک‌شدن به صفر، در محیطی با هزینه‌ی فرصتِ بالا، مشاهده
می‌کند درمی‌یابد که ارزشِ واقعیِ پرداخت‌های آتیِ نقطه‌ی عطف در حالِ
فرسایش است درحالی‌که ارزشِ تسویه‌ی کارِ انجام‌شده به‌لحاظِ اسمی ثابت
می‌ماند ---و این خود انگیزه‌ای طبیعی برای تخصیصِ مجدد ایجاد می‌کند که
مدل از طریقِ تعاملِ ضریبِ تنزیل و سازوکارهای تسویه‌ی فسخ، بدونِ نیاز
به مدل‌سازیِ صریحِ تورم، آن را ثبت می‌کند~\citep{Dixit1994,
FIDIC2017}.

% ------------------------------------------------------------
\subsection{یادگیری تقویتی برای تخصیص ترتیبی منابع}
\label{subsec:lr_rl}
% ------------------------------------------------------------

یادگیریِ تقویتی چارچوبی اصولی برای تصمیم‌گیریِ ترتیبی تحتِ عدم‌قطعیت
فراهم می‌کند~\citep{Sutton2018}. صورت‌بندیِ \lr{MDP} ---حالت، عمل،
گذار، پاداش--- به‌طورِ طبیعی بر مسئله‌ی تخصیصِ بودجه‌ی سبد منطبق
می‌شود: مدیرِ سبد به‌صورتِ سیاستی مدل می‌شود که حالتِ جاری را مشاهده
می‌کند، تخصیصی را متعهد می‌شود، و پاداشی از جریانِ نقدیِ خالصِ
تنزیل‌شده دریافت می‌کند، در حالی‌که محیط به‌صورتِ تصادفی و بر اساسِ
بازدهی‌های محقق‌شده گذار می‌کند.

\lr{RL} برای زمان‌بندیِ پروژه~\citep{Peng2021, Liu2023}، مسائلِ
زمان‌بندیِ پروژه‌ی مقیدشده به منابع~\citep{Tassel2023}، و مدیریتِ
موجودیِ زنجیره‌ی تأمین~\citep{Giannoccaro2002} به‌کار رفته است. در هر
یک از این حوزه‌ها، \lr{RL} در مواردی که فضای حالت بزرگ است، دینامیکِ
گذار تصادفی است، مسئله در طولِ افقِ تصمیم‌گیری ناایستا است، و طولِ
افق در آغازِ برنامه‌ریزی نامعلوم است، نسبت به روش‌های کلاسیکِ تحقیق
در عملیات برتری نشان داده است. این‌ها دقیقاً همان ویژگی‌هایی هستند که
مسئله‌ی تخصیصِ بودجه‌ی سبد دارد.

رابطِ برنامه‌نویسیِ \lr{Gymnasium}~\citep{Towers2024} به‌عنوانِ رابطِ
استانداردِ \textit{de facto} برای توسعه‌ی محیط‌های \lr{RL} عمل
می‌کند و ادغامِ مستقیم با کتابخانه‌های مدرنِ بهینه‌سازیِ سیاست را
ممکن می‌سازد. تصادفی‌سازیِ حوزه~\citep{Tobin2017} ---تصادفی‌سازیِ
پارامترهای محیط در طولِ قسمت‌های آموزشی--- به‌صورتِ بومی از طریقِ
دیکشنریِ پیکربندی پشتیبانی می‌شود و به یک سیاستِ آموزش‌دیده امکان
می‌دهد بدونِ نیاز به آموزشِ مجدد، بر رویِ ساختارهای قراردادی و
انواعِ پروژه‌ی متنوع تعمیم یابد.

کاربردِ مستقیمِ \lr{RL} برای تخصیصِ بودجه در سطحِ سبد، تحتِ
سازوکارهای جریانِ نقدیِ قراردادی، همچنان در ادبیات غایب است. پژوهشِ
حاضر نخستین صورت‌بندیِ \lr{MDP} از این مسئله و نخستین نشانِ تجربی را
ارائه می‌دهد که اکتشافی‌های قطعیِ حریصانه، با رشدِ پیچیدگیِ سبد،
به‌طورِ نظام‌مند نسبت به بهینه‌ی اطلاعاتِ‌کامل ضعیف‌تر عمل می‌کنند،
و بدین‌ترتیب هدفی بهبودِ کمّی‌شده برایِ سیاست‌های یادگیری‌شده‌ی آینده
تعیین می‌کند.

% ============================================================
%  نتیجه‌گیری
% ============================================================
\section{نتیجه‌گیری}
\label{subsec:lr_synthesis}

جدول~\ref{tab:litmap}، آثارِ نماینده را در چهار جریانِ ادبیات در برابرِ
ابعادِ مدل‌سازیِ کلیدیِ چارچوبِ حاضر نگاشت می‌کند، همان‌گونه که توسطِ
خطِ‌لوله‌ی سنتزِ خودکارِ شرح‌داده‌شده در بخش~\ref{subsubsec:prisma}
تولید شده است. سه شکافِ ساختاری از این نگاشت پدیدار می‌شود.

\textbf{شکافِ ۱ --- انتخاب در برابرِ حاکمیتِ عملیاتی.}\;
گذار از \emph{انتخابِ} سبد به \emph{حاکمیتِ عملیاتیِ} سبد، به‌عنوانِ
مسئله‌ای از تصمیم‌گیریِ ترتیبی مدل نشده است. ادبیاتِ انتخاب،
مدیریتِ پس‌از‌انتخاب را نادیده می‌گیرد؛ ادبیاتِ زمان‌بندیِ پروژه نیز
وابستگی‌های متقابلِ سبد و قیدِ سرمایه‌ی مشترک را نادیده می‌گیرد. هیچ
اثرِ موجودی هر دو را هم‌زمان پوشش نمی‌دهد.

\textbf{شکافِ ۲ --- انتشارِ عدم‌قطعیت از سطحِ پروژه به سطحِ سبد.}\;
مدل‌های تصادفیِ موجود از عدم‌قطعیتِ جریانِ نقدیِ پروژه، در سطحِ
پروژه‌ی منفرد عمل می‌کنند و عدم‌قطعیت را به معماریِ تصمیم‌گیریِ
سطحِ سبد تحتِ حاکمیتِ قیدِ بودجه‌ی مشترک منتقل نمی‌کنند. حلقه‌ی
بازخوردِ میانِ مسیرِ بازدهیِ تصادفیِ یک پروژه و حقِ دسترسیِ آن به
تأمینِ مالی در سطحِ سبد، همچنان بدونِ مدل باقی مانده است.

\textbf{شکافِ ۳ --- غیابِ سازوکارهای درون‌زای فسخ.}\;
مدل‌سازیِ مشترکِ یک شرطِ نقضِ تأخیرِ پیشرفت ردیابی‌شده در برابرِ
تحملِ قراردادی، یک شمارنده‌ی دوره‌ی مهلتِ کاهشی که فرسایشِ دوره‌ی
مهلت را اداره می‌کند، یک شرطِ نقضِ مهلتِ سختِ مستقل که فسخِ فوری را
صرفِ‌نظر از وضعیتِ شمارنده‌ی دوره‌ی مهلت فعال می‌کند، و جریانِ نقدیِ
تسویه‌ی فسخ که ارزشِ کسب‌شده را در برابرِ پرداخت‌های تجمعیِ دریافتی
تطبیق می‌دهد، تاکنون در چارچوبی از تصمیم‌گیریِ ترتیبی در سطحِ سبد که
در آن تخصیصِ بودجه تنها متغیرِ کنترلی باشد، جاسازی نشده است.

\medskip
این شکاف‌ها خطاهایی منفرد نیستند؛ آن‌ها پیشرفتی ناقص در مسیرِ
دستیابی به سامانه‌ای عملیاتیِ پشتیبانِ تصمیم در سطحِ سبد را بازتاب
می‌دهند. شکل~\ref{fig:maturity} این پیشرفت را در قالبِ سه گامِ
متوالی جای می‌دهد.

\textbf{گامِ ۱}---زمینه‌سازیِ حوزه---به‌خوبی توسطِ ادبیاتِ مرورشده در
این فصل تثبیت شده است: سازوکارهای قراردادی، سیگنال‌های عملکردِ
\lr{EVM}، پروفیل‌های مخارجِ منحنیِ \lr{S}، و مدل‌های انتخابِ سبد،
در مجموع ورودی‌ها و شرایطِ مرزی‌ای را فراهم می‌کنند که هر مدلِ
حاکمیتِ عملیاتی به آن‌ها نیاز دارد.

\textbf{گامِ ۲}---صورت‌بندیِ مسئله‌ای قابلِ‌حل---شکافی است که این
پژوهش آن را پر می‌کند. دستاوردِ این پژوهش، محیطی از نوعِ \lr{MDP}
است، تأییدشده و پیکربندی‌پذیر، که چرخه‌ی کاملِ جریانِ نقدیِ
\lr{FIDIC} و سازوکارِ فسخِ دوشرطی را رمزگذاری می‌کند و به‌صورتِ
بسته‌ای سازگار با \lr{Gymnasium} به‌همراه دفترِ‌کلِ حسابرسیِ
\lr{SQLite}ِ کامل منتشر شده است. اکتشافی‌های قطعیِ حریصانه در برابرِ
بهینه‌ی اطلاعاتِ‌کامل محک زده شده‌اند و شکافِ بهینگیِ کمّی‌شده‌ای به
دست آمده که به‌طورِ یکنواخت با پیچیدگیِ سبد رشد می‌کند و هدفِ عملکردی
برای عامل‌های یادگیری‌شده‌ی آینده را تعیین می‌کند.

\textbf{گامِ ۳}---استقرارِ اعتبارسنجی‌شده، شاملِ آموزشِ عاملِ \lr{RL}
و ادغام با سامانه‌های اطلاعاتِ مدیریتِ پروژه‌ی رایج مانندِ
\lr{Oracle Primavera~P6}--- همچنان کارِ آینده باقی می‌ماند و مشروط به
محیطی است که گامِ ۲ اکنون فراهم می‌کند.

\begin{figure}[htbp]
    \centering
    % figures/maturity.tex
% TikZ only — no figure environment, no caption, no label.
% Include in section file as:
%   % figures/maturity.tex
% TikZ only — no figure environment, no caption, no label.
% Include in section file as:
%   % figures/maturity.tex
% TikZ only — no figure environment, no caption, no label.
% Include in section file as:
%   \input{figures/maturity.tex}
\begin{tikzpicture}[
    node distance  = 0cm,
    stepbox/.style = {rectangle, draw, rounded corners=3pt,
                      minimum width=3.6cm, minimum height=2.8cm,
                      align=center, font=\small, text width=3.2cm},
    arrow/.style   = {->, >=stealth, thick},
  ]

  \node[stepbox, fill=gray!15] (s1) at (0,0)
    {\textbf{Step 1}\\[4pt]
     Domain Grounding\\[6pt]
     \scriptsize
     Contract mechanics\\
     EVM signals\\
     S-curve profiles\\
     Portfolio selection};

  \node[stepbox, fill=gray!8, right=0.6cm of s1] (s2)
    {\textbf{Step 2}\\[1pt]
     \emph{(this work)}\\[4pt]
     Problem Formulation\\[6pt]
     \scriptsize
     MDP environment\\
     FIDIC cash-flow\\
     Termination mechanic\\
     Gymnasium API};

  \node[stepbox, draw=gray, dashed, right=0.6cm of s2] (s3)
    {\textbf{Step 3}\\[4pt]
     Validated Deployment\\[6pt]
     \scriptsize
     RL agent training\\
     OR benchmark\\
     Primavera~P6 integration\\[4pt]
     \emph{(future work)}};

  \draw[arrow]                    (s1) -- (s2);
  \draw[arrow, dashed, color=gray] (s2) -- (s3);

\end{tikzpicture}
\begin{tikzpicture}[
    node distance  = 0cm,
    stepbox/.style = {rectangle, draw, rounded corners=3pt,
                      minimum width=3.6cm, minimum height=2.8cm,
                      align=center, font=\small, text width=3.2cm},
    arrow/.style   = {->, >=stealth, thick},
  ]

  \node[stepbox, fill=gray!15] (s1) at (0,0)
    {\textbf{Step 1}\\[4pt]
     Domain Grounding\\[6pt]
     \scriptsize
     Contract mechanics\\
     EVM signals\\
     S-curve profiles\\
     Portfolio selection};

  \node[stepbox, fill=gray!8, right=0.6cm of s1] (s2)
    {\textbf{Step 2}\\[1pt]
     \emph{(this work)}\\[4pt]
     Problem Formulation\\[6pt]
     \scriptsize
     MDP environment\\
     FIDIC cash-flow\\
     Termination mechanic\\
     Gymnasium API};

  \node[stepbox, draw=gray, dashed, right=0.6cm of s2] (s3)
    {\textbf{Step 3}\\[4pt]
     Validated Deployment\\[6pt]
     \scriptsize
     RL agent training\\
     OR benchmark\\
     Primavera~P6 integration\\[4pt]
     \emph{(future work)}};

  \draw[arrow]                    (s1) -- (s2);
  \draw[arrow, dashed, color=gray] (s2) -- (s3);

\end{tikzpicture}
\begin{tikzpicture}[
    node distance  = 0cm,
    stepbox/.style = {rectangle, draw, rounded corners=3pt,
                      minimum width=3.6cm, minimum height=2.8cm,
                      align=center, font=\small, text width=3.2cm},
    arrow/.style   = {->, >=stealth, thick},
  ]

  \node[stepbox, fill=gray!15] (s1) at (0,0)
    {\textbf{Step 1}\\[4pt]
     Domain Grounding\\[6pt]
     \scriptsize
     Contract mechanics\\
     EVM signals\\
     S-curve profiles\\
     Portfolio selection};

  \node[stepbox, fill=gray!8, right=0.6cm of s1] (s2)
    {\textbf{Step 2}\\[1pt]
     \emph{(this work)}\\[4pt]
     Problem Formulation\\[6pt]
     \scriptsize
     MDP environment\\
     FIDIC cash-flow\\
     Termination mechanic\\
     Gymnasium API};

  \node[stepbox, draw=gray, dashed, right=0.6cm of s2] (s3)
    {\textbf{Step 3}\\[4pt]
     Validated Deployment\\[6pt]
     \scriptsize
     RL agent training\\
     OR benchmark\\
     Primavera~P6 integration\\[4pt]
     \emph{(future work)}};

  \draw[arrow]                    (s1) -- (s2);
  \draw[arrow, dashed, color=gray] (s2) -- (s3);

\end{tikzpicture}
    \caption{روندِ بلوغ به‌سویِ سامانه‌ای عملیاتیِ پشتیبانِ تصمیم در سطحِ
      سبد. گامِ ۲، که توسطِ پژوهشِ حاضر تکمیل شده، گامِ ۳ را ممکن
      می‌سازد. حاشیه و فلشِ نقطه‌چین، کارِ آینده را نشان می‌دهند.}
    \label{fig:maturity}
\end{figure}

\medskip
پژوهشِ حاضر هر سه شکاف را پوشش می‌دهد و گامِ ۲ از این روند را تکمیل
می‌کند. محیطِ \lr{MDP} سازگار با \lr{Gymnasium} ---با سازوکارهای
قراردادیِ پارامتری‌پذیر، بازدهیِ تصادفیِ بودجه‌به‌پیشرفت به‌عنوانِ
تنها منبعِ عدم‌قطعیت، و سازوکارِ فسخِ دوشرطیِ درون‌زا--- دستاوردِ
روش‌شناختیِ اصلیِ این پژوهش را تشکیل می‌دهد. اکتشافی‌های قطعیِ
حریصانه در برابرِ بهینه‌ی اطلاعاتِ‌کامل در بخش~\ref{sec:evaluation}
ارزیابی می‌شوند و شکافِ بهینگیِ کمّی‌شده‌ای ارائه می‌دهند که به‌طورِ
یکنواخت با پیچیدگیِ سبد رشد می‌کند و هدفِ بهبود برای عامل‌های
\lr{RL} آینده را تعریف می‌کند.

\begin{table}[htbp]
    \centering
    \caption{آثارِ نماینده، نگاشت‌شده در برابرِ ابعادِ مدل‌سازیِ
      چارچوبِ حاضر. \textit{نوعِ تصمیم}: الف\,=\,انتخاب،
      ز\,=\,زمان‌بندی، ب\,=\,بودجه‌بندی. \checkmark~=~به‌طورِ کامل
      پوشش‌داده‌شده؛\; $\circ$~=~تاحدی پوشش‌داده‌شده؛\;
      {--}~=~پوشش‌داده‌نشده.}
    \label{tab:litmap}
    % tables/litmap.tex
% tabular only — no table environment, no caption, no label.
% Include in section file as:
%   % tables/litmap.tex
% tabular only — no table environment, no caption, no label.
% Include in section file as:
%   % tables/litmap.tex
% tabular only — no table environment, no caption, no label.
% Include in section file as:
%   \input{tables/litmap.tex}
\setlength{\tabcolsep}{6pt}
\renewcommand{\arraystretch}{1.35}
\rowcolors{3}{rowodd}{white}
\resizebox{\textwidth}{!}{%
\begin{tabular}{>{\raggedright\arraybackslash}p{3.8cm}
                c
                c
                >{\columncolor{rarebg}}c
                c
                >{\columncolor{rarebg}}c
                >{\columncolor{rarebg}}c}
  \toprule
  \rowcolor{headbg}
    \color{headfg}\textbf{Work}
    & \color{headfg}\textbf{Portfolio}
    & \color{headfg}\textbf{Decision}
    & \cellcolor{headbg}\color{headfg}\textbf{Sequential}
    & \color{headfg}\textbf{Stochastic}
    & \cellcolor{headbg}\color{headfg}\textbf{Contract}
    & \cellcolor{headbg}\color{headfg}\textbf{Termination} \\
  \rowcolor{headbg}
    \color{headfg}\textbf{}
    & \color{headfg}\textbf{level}
    & \color{headfg}\textbf{type}
    & \cellcolor{headbg}\color{headfg}\textbf{decisions}
    & \color{headfg}\textbf{cash flows}
    & \cellcolor{headbg}\color{headfg}\textbf{mechanics}
    & \cellcolor{headbg}\color{headfg}\textbf{settlement} \\
  \midrule
  \citet{Yamakawa2019}        & \textcolor{symgreen}{\checkmark} & \textsc{\small S}    & \textcolor{symgrey}{--}          & \textcolor{symamber}{$\circ$}    & \textcolor{symgrey}{--}          & \textcolor{symgrey}{--}          \\
  \citet{KOSZTYAN2022}        & \textcolor{symgreen}{\checkmark} & \textsc{\small S}    & \textcolor{symgrey}{--}          & \textcolor{symamber}{$\circ$}    & \textcolor{symgrey}{--}          & \textcolor{symgrey}{--}          \\
  \citet{touran2006mathematical}  & \textcolor{symgreen}{\checkmark} & \textsc{\small Sc/B} & \textcolor{symgrey}{--}          & \textcolor{symamber}{$\circ$}    & \textcolor{symgrey}{--}          & \textcolor{symgrey}{--}          \\
  \citet{touran2010probabilistic} & \textcolor{symgreen}{\checkmark} & \textsc{\small Sc/B} & \textcolor{symgrey}{--}          & \textcolor{symamber}{$\circ$}    & \textcolor{symgrey}{--}          & \textcolor{symgrey}{--}          \\
  \citet{GHAPANCHI2012791}    & \textcolor{symgreen}{\checkmark} & \textsc{\small S}    & \textcolor{symgrey}{--}          & \textcolor{symamber}{$\circ$}    & \textcolor{symgrey}{--}          & \textcolor{symgrey}{--}          \\
  \citet{pourahmadi2015scenario}  & \textcolor{symgreen}{\checkmark} & \textsc{\small S}    & \textcolor{symgrey}{--}          & \textcolor{symamber}{$\circ$}    & \textcolor{symgrey}{--}          & \textcolor{symgrey}{--}          \\
  \citet{elsabek2018coordination} & \textcolor{symgrey}{--}          & \textsc{\small B}    & \textcolor{symgrey}{--}          & \textcolor{symgrey}{--}          & \textcolor{symgreen}{\checkmark} & \textcolor{symgrey}{--}          \\
  \citet{beseiso2021}         & \textcolor{symgreen}{\checkmark} & \textsc{\small S}    & \textcolor{symgrey}{--}          & \textcolor{symamber}{$\circ$}    & \textcolor{symgrey}{--}          & \textcolor{symgrey}{--}          \\
  \citet{karimi2021scenario}  & \textcolor{symgreen}{\checkmark} & \textsc{\small S}    & \textcolor{symgrey}{--}          & \textcolor{symamber}{$\circ$}    & \textcolor{symgrey}{--}          & \textcolor{symgrey}{--}          \\
  \citet{bahroun2023artificial}   & \textcolor{symgrey}{--}          & \textsc{\small Sc}   & \textcolor{symgreen}{\checkmark} & \textcolor{symgrey}{--}          & \textcolor{symgrey}{--}          & \textcolor{symgrey}{--}          \\
  \citet{RANJBAR2022116134}   & \textcolor{symgreen}{\checkmark} & \textsc{\small S/Sc} & \textcolor{symgrey}{--}          & \textcolor{symgrey}{--}          & \textcolor{symamber}{$\circ$}    & \textcolor{symgrey}{--}          \\
  \citet{szwarcfiter2023balancing}& \textcolor{symgreen}{\checkmark} & \textsc{\small Sc}   & \textcolor{symgreen}{\checkmark} & \textcolor{symamber}{$\circ$}    & \textcolor{symgrey}{--}          & \textcolor{symgrey}{--}          \\
  \citet{asghari2023}         & \textcolor{symgrey}{--}          & \textcolor{symgrey}{--} & \textcolor{symgreen}{\checkmark} & \textcolor{symamber}{$\circ$} & \textcolor{symgrey}{--}          & \textcolor{symgrey}{--}          \\
  \citet{mazhari2024}         & \textcolor{symgreen}{\checkmark} & \textsc{\small S/B}  & \textcolor{symgreen}{\checkmark} & \textcolor{symgrey}{--}          & \textcolor{symgrey}{--}          & \textcolor{symgrey}{--}          \\
  \citet{JIANG2024101016}     & \textcolor{symgreen}{\checkmark} & \textcolor{symgrey}{--} & \textcolor{symgreen}{\checkmark} & \textcolor{symamber}{$\circ$} & \textcolor{symgrey}{--}          & \textcolor{symgrey}{--}          \\
  \citet{Nabot2025}           & \textcolor{symgreen}{\checkmark} & \textsc{\small Sc}   & \textcolor{symgreen}{\checkmark} & \textcolor{symgrey}{--}          & \textcolor{symgreen}{\checkmark} & \textcolor{symgreen}{\checkmark} \\
  \citet{Darvish2025}         & \textcolor{symgreen}{\checkmark} & \textsc{\small B}    & \textcolor{symgreen}{\checkmark} & \textcolor{symamber}{$\circ$}    & \textcolor{symamber}{$\circ$}    & \textcolor{symgreen}{\checkmark} \\
  \citet{Gholami2025}         & \textcolor{symgreen}{\checkmark} & \textsc{\small Sc}   & \textcolor{symgreen}{\checkmark} & \textcolor{symamber}{$\circ$}    & \textcolor{symgrey}{--}          & \textcolor{symgreen}{\checkmark} \\
  \citet{Kehinde2025}         & \textcolor{symgreen}{\checkmark} & \textsc{\small B}    & \textcolor{symgreen}{\checkmark} & \textcolor{symgreen}{\checkmark} & \textcolor{symgrey}{--}          & \textcolor{symgrey}{--}          \\
  \citet{Ajibade2025}         & \textcolor{symgreen}{\checkmark} & \textsc{\small Sc}   & \textcolor{symgreen}{\checkmark} & \textcolor{symamber}{$\circ$}    & \textcolor{symgrey}{--}          & \textcolor{symgrey}{--}          \\
  \citet{Huddar2025}          & \textcolor{symgreen}{\checkmark} & \textsc{\small B}    & \textcolor{symgreen}{\checkmark} & \textcolor{symamber}{$\circ$}    & \textcolor{symgrey}{--}          & \textcolor{symgrey}{--}          \\
  \citet{Akbari2026}          & \textcolor{symgreen}{\checkmark} & \textsc{\small B}    & \textcolor{symgreen}{\checkmark} & \textcolor{symamber}{$\circ$}    & \textcolor{symgreen}{\checkmark} & \textcolor{symgreen}{\checkmark} \\
  \citet{ebnerasoul2023}      & \textcolor{symgreen}{\checkmark} & \textsc{\small S}    & \textcolor{symgrey}{--}          & \textcolor{symamber}{$\circ$}    & \textcolor{symgrey}{--}          & \textcolor{symgrey}{--}          \\
  \citet{noroozi2016probabilistic}& \textcolor{symgreen}{\checkmark} & \textsc{\small B}  & \textcolor{symgrey}{--}          & \textcolor{symamber}{$\circ$}    & \textcolor{symgrey}{--}          & \textcolor{symgrey}{--}          \\
  \citet{JIANG2023104817}     & \textcolor{symgrey}{--}          & \textsc{\small B}    & \textcolor{symgreen}{\checkmark} & \textcolor{symgrey}{--}          & \textcolor{symgrey}{--}          & \textcolor{symgrey}{--}          \\
  \cmidrule{1-7}
  \rowcolor{pwbg}
    \color{headfg}\textbf{Present work}
    & \cellcolor{pwbg}\color{headfg}\checkmark
    & \color{headfg}\textbf{B}
    & \cellcolor{pwbg}\color{headfg}\checkmark
    & \color{headfg}\checkmark
    & \cellcolor{pwbg}\color{headfg}\checkmark
    & \cellcolor{pwbg}\color{headfg}\checkmark \\
  \bottomrule
\end{tabular}%
}
\setlength{\tabcolsep}{6pt}
\renewcommand{\arraystretch}{1.35}
\rowcolors{3}{rowodd}{white}
\resizebox{\textwidth}{!}{%
\begin{tabular}{>{\raggedright\arraybackslash}p{3.8cm}
                c
                c
                >{\columncolor{rarebg}}c
                c
                >{\columncolor{rarebg}}c
                >{\columncolor{rarebg}}c}
  \toprule
  \rowcolor{headbg}
    \color{headfg}\textbf{Work}
    & \color{headfg}\textbf{Portfolio}
    & \color{headfg}\textbf{Decision}
    & \cellcolor{headbg}\color{headfg}\textbf{Sequential}
    & \color{headfg}\textbf{Stochastic}
    & \cellcolor{headbg}\color{headfg}\textbf{Contract}
    & \cellcolor{headbg}\color{headfg}\textbf{Termination} \\
  \rowcolor{headbg}
    \color{headfg}\textbf{}
    & \color{headfg}\textbf{level}
    & \color{headfg}\textbf{type}
    & \cellcolor{headbg}\color{headfg}\textbf{decisions}
    & \color{headfg}\textbf{cash flows}
    & \cellcolor{headbg}\color{headfg}\textbf{mechanics}
    & \cellcolor{headbg}\color{headfg}\textbf{settlement} \\
  \midrule
  \citet{Yamakawa2019}        & \textcolor{symgreen}{\checkmark} & \textsc{\small S}    & \textcolor{symgrey}{--}          & \textcolor{symamber}{$\circ$}    & \textcolor{symgrey}{--}          & \textcolor{symgrey}{--}          \\
  \citet{KOSZTYAN2022}        & \textcolor{symgreen}{\checkmark} & \textsc{\small S}    & \textcolor{symgrey}{--}          & \textcolor{symamber}{$\circ$}    & \textcolor{symgrey}{--}          & \textcolor{symgrey}{--}          \\
  \citet{touran2006mathematical}  & \textcolor{symgreen}{\checkmark} & \textsc{\small Sc/B} & \textcolor{symgrey}{--}          & \textcolor{symamber}{$\circ$}    & \textcolor{symgrey}{--}          & \textcolor{symgrey}{--}          \\
  \citet{touran2010probabilistic} & \textcolor{symgreen}{\checkmark} & \textsc{\small Sc/B} & \textcolor{symgrey}{--}          & \textcolor{symamber}{$\circ$}    & \textcolor{symgrey}{--}          & \textcolor{symgrey}{--}          \\
  \citet{GHAPANCHI2012791}    & \textcolor{symgreen}{\checkmark} & \textsc{\small S}    & \textcolor{symgrey}{--}          & \textcolor{symamber}{$\circ$}    & \textcolor{symgrey}{--}          & \textcolor{symgrey}{--}          \\
  \citet{pourahmadi2015scenario}  & \textcolor{symgreen}{\checkmark} & \textsc{\small S}    & \textcolor{symgrey}{--}          & \textcolor{symamber}{$\circ$}    & \textcolor{symgrey}{--}          & \textcolor{symgrey}{--}          \\
  \citet{elsabek2018coordination} & \textcolor{symgrey}{--}          & \textsc{\small B}    & \textcolor{symgrey}{--}          & \textcolor{symgrey}{--}          & \textcolor{symgreen}{\checkmark} & \textcolor{symgrey}{--}          \\
  \citet{beseiso2021}         & \textcolor{symgreen}{\checkmark} & \textsc{\small S}    & \textcolor{symgrey}{--}          & \textcolor{symamber}{$\circ$}    & \textcolor{symgrey}{--}          & \textcolor{symgrey}{--}          \\
  \citet{karimi2021scenario}  & \textcolor{symgreen}{\checkmark} & \textsc{\small S}    & \textcolor{symgrey}{--}          & \textcolor{symamber}{$\circ$}    & \textcolor{symgrey}{--}          & \textcolor{symgrey}{--}          \\
  \citet{bahroun2023artificial}   & \textcolor{symgrey}{--}          & \textsc{\small Sc}   & \textcolor{symgreen}{\checkmark} & \textcolor{symgrey}{--}          & \textcolor{symgrey}{--}          & \textcolor{symgrey}{--}          \\
  \citet{RANJBAR2022116134}   & \textcolor{symgreen}{\checkmark} & \textsc{\small S/Sc} & \textcolor{symgrey}{--}          & \textcolor{symgrey}{--}          & \textcolor{symamber}{$\circ$}    & \textcolor{symgrey}{--}          \\
  \citet{szwarcfiter2023balancing}& \textcolor{symgreen}{\checkmark} & \textsc{\small Sc}   & \textcolor{symgreen}{\checkmark} & \textcolor{symamber}{$\circ$}    & \textcolor{symgrey}{--}          & \textcolor{symgrey}{--}          \\
  \citet{asghari2023}         & \textcolor{symgrey}{--}          & \textcolor{symgrey}{--} & \textcolor{symgreen}{\checkmark} & \textcolor{symamber}{$\circ$} & \textcolor{symgrey}{--}          & \textcolor{symgrey}{--}          \\
  \citet{mazhari2024}         & \textcolor{symgreen}{\checkmark} & \textsc{\small S/B}  & \textcolor{symgreen}{\checkmark} & \textcolor{symgrey}{--}          & \textcolor{symgrey}{--}          & \textcolor{symgrey}{--}          \\
  \citet{JIANG2024101016}     & \textcolor{symgreen}{\checkmark} & \textcolor{symgrey}{--} & \textcolor{symgreen}{\checkmark} & \textcolor{symamber}{$\circ$} & \textcolor{symgrey}{--}          & \textcolor{symgrey}{--}          \\
  \citet{Nabot2025}           & \textcolor{symgreen}{\checkmark} & \textsc{\small Sc}   & \textcolor{symgreen}{\checkmark} & \textcolor{symgrey}{--}          & \textcolor{symgreen}{\checkmark} & \textcolor{symgreen}{\checkmark} \\
  \citet{Darvish2025}         & \textcolor{symgreen}{\checkmark} & \textsc{\small B}    & \textcolor{symgreen}{\checkmark} & \textcolor{symamber}{$\circ$}    & \textcolor{symamber}{$\circ$}    & \textcolor{symgreen}{\checkmark} \\
  \citet{Gholami2025}         & \textcolor{symgreen}{\checkmark} & \textsc{\small Sc}   & \textcolor{symgreen}{\checkmark} & \textcolor{symamber}{$\circ$}    & \textcolor{symgrey}{--}          & \textcolor{symgreen}{\checkmark} \\
  \citet{Kehinde2025}         & \textcolor{symgreen}{\checkmark} & \textsc{\small B}    & \textcolor{symgreen}{\checkmark} & \textcolor{symgreen}{\checkmark} & \textcolor{symgrey}{--}          & \textcolor{symgrey}{--}          \\
  \citet{Ajibade2025}         & \textcolor{symgreen}{\checkmark} & \textsc{\small Sc}   & \textcolor{symgreen}{\checkmark} & \textcolor{symamber}{$\circ$}    & \textcolor{symgrey}{--}          & \textcolor{symgrey}{--}          \\
  \citet{Huddar2025}          & \textcolor{symgreen}{\checkmark} & \textsc{\small B}    & \textcolor{symgreen}{\checkmark} & \textcolor{symamber}{$\circ$}    & \textcolor{symgrey}{--}          & \textcolor{symgrey}{--}          \\
  \citet{Akbari2026}          & \textcolor{symgreen}{\checkmark} & \textsc{\small B}    & \textcolor{symgreen}{\checkmark} & \textcolor{symamber}{$\circ$}    & \textcolor{symgreen}{\checkmark} & \textcolor{symgreen}{\checkmark} \\
  \citet{ebnerasoul2023}      & \textcolor{symgreen}{\checkmark} & \textsc{\small S}    & \textcolor{symgrey}{--}          & \textcolor{symamber}{$\circ$}    & \textcolor{symgrey}{--}          & \textcolor{symgrey}{--}          \\
  \citet{noroozi2016probabilistic}& \textcolor{symgreen}{\checkmark} & \textsc{\small B}  & \textcolor{symgrey}{--}          & \textcolor{symamber}{$\circ$}    & \textcolor{symgrey}{--}          & \textcolor{symgrey}{--}          \\
  \citet{JIANG2023104817}     & \textcolor{symgrey}{--}          & \textsc{\small B}    & \textcolor{symgreen}{\checkmark} & \textcolor{symgrey}{--}          & \textcolor{symgrey}{--}          & \textcolor{symgrey}{--}          \\
  \cmidrule{1-7}
  \rowcolor{pwbg}
    \color{headfg}\textbf{Present work}
    & \cellcolor{pwbg}\color{headfg}\checkmark
    & \color{headfg}\textbf{B}
    & \cellcolor{pwbg}\color{headfg}\checkmark
    & \color{headfg}\checkmark
    & \cellcolor{pwbg}\color{headfg}\checkmark
    & \cellcolor{pwbg}\color{headfg}\checkmark \\
  \bottomrule
\end{tabular}%
}
\setlength{\tabcolsep}{6pt}
\renewcommand{\arraystretch}{1.35}
\rowcolors{3}{rowodd}{white}
\resizebox{\textwidth}{!}{%
\begin{tabular}{>{\raggedright\arraybackslash}p{3.8cm}
                c
                c
                >{\columncolor{rarebg}}c
                c
                >{\columncolor{rarebg}}c
                >{\columncolor{rarebg}}c}
  \toprule
  \rowcolor{headbg}
    \color{headfg}\textbf{Work}
    & \color{headfg}\textbf{Portfolio}
    & \color{headfg}\textbf{Decision}
    & \cellcolor{headbg}\color{headfg}\textbf{Sequential}
    & \color{headfg}\textbf{Stochastic}
    & \cellcolor{headbg}\color{headfg}\textbf{Contract}
    & \cellcolor{headbg}\color{headfg}\textbf{Termination} \\
  \rowcolor{headbg}
    \color{headfg}\textbf{}
    & \color{headfg}\textbf{level}
    & \color{headfg}\textbf{type}
    & \cellcolor{headbg}\color{headfg}\textbf{decisions}
    & \color{headfg}\textbf{cash flows}
    & \cellcolor{headbg}\color{headfg}\textbf{mechanics}
    & \cellcolor{headbg}\color{headfg}\textbf{settlement} \\
  \midrule
  \citet{Yamakawa2019}        & \textcolor{symgreen}{\checkmark} & \textsc{\small S}    & \textcolor{symgrey}{--}          & \textcolor{symamber}{$\circ$}    & \textcolor{symgrey}{--}          & \textcolor{symgrey}{--}          \\
  \citet{KOSZTYAN2022}        & \textcolor{symgreen}{\checkmark} & \textsc{\small S}    & \textcolor{symgrey}{--}          & \textcolor{symamber}{$\circ$}    & \textcolor{symgrey}{--}          & \textcolor{symgrey}{--}          \\
  \citet{touran2006mathematical}  & \textcolor{symgreen}{\checkmark} & \textsc{\small Sc/B} & \textcolor{symgrey}{--}          & \textcolor{symamber}{$\circ$}    & \textcolor{symgrey}{--}          & \textcolor{symgrey}{--}          \\
  \citet{touran2010probabilistic} & \textcolor{symgreen}{\checkmark} & \textsc{\small Sc/B} & \textcolor{symgrey}{--}          & \textcolor{symamber}{$\circ$}    & \textcolor{symgrey}{--}          & \textcolor{symgrey}{--}          \\
  \citet{GHAPANCHI2012791}    & \textcolor{symgreen}{\checkmark} & \textsc{\small S}    & \textcolor{symgrey}{--}          & \textcolor{symamber}{$\circ$}    & \textcolor{symgrey}{--}          & \textcolor{symgrey}{--}          \\
  \citet{pourahmadi2015scenario}  & \textcolor{symgreen}{\checkmark} & \textsc{\small S}    & \textcolor{symgrey}{--}          & \textcolor{symamber}{$\circ$}    & \textcolor{symgrey}{--}          & \textcolor{symgrey}{--}          \\
  \citet{elsabek2018coordination} & \textcolor{symgrey}{--}          & \textsc{\small B}    & \textcolor{symgrey}{--}          & \textcolor{symgrey}{--}          & \textcolor{symgreen}{\checkmark} & \textcolor{symgrey}{--}          \\
  \citet{beseiso2021}         & \textcolor{symgreen}{\checkmark} & \textsc{\small S}    & \textcolor{symgrey}{--}          & \textcolor{symamber}{$\circ$}    & \textcolor{symgrey}{--}          & \textcolor{symgrey}{--}          \\
  \citet{karimi2021scenario}  & \textcolor{symgreen}{\checkmark} & \textsc{\small S}    & \textcolor{symgrey}{--}          & \textcolor{symamber}{$\circ$}    & \textcolor{symgrey}{--}          & \textcolor{symgrey}{--}          \\
  \citet{bahroun2023artificial}   & \textcolor{symgrey}{--}          & \textsc{\small Sc}   & \textcolor{symgreen}{\checkmark} & \textcolor{symgrey}{--}          & \textcolor{symgrey}{--}          & \textcolor{symgrey}{--}          \\
  \citet{RANJBAR2022116134}   & \textcolor{symgreen}{\checkmark} & \textsc{\small S/Sc} & \textcolor{symgrey}{--}          & \textcolor{symgrey}{--}          & \textcolor{symamber}{$\circ$}    & \textcolor{symgrey}{--}          \\
  \citet{szwarcfiter2023balancing}& \textcolor{symgreen}{\checkmark} & \textsc{\small Sc}   & \textcolor{symgreen}{\checkmark} & \textcolor{symamber}{$\circ$}    & \textcolor{symgrey}{--}          & \textcolor{symgrey}{--}          \\
  \citet{asghari2023}         & \textcolor{symgrey}{--}          & \textcolor{symgrey}{--} & \textcolor{symgreen}{\checkmark} & \textcolor{symamber}{$\circ$} & \textcolor{symgrey}{--}          & \textcolor{symgrey}{--}          \\
  \citet{mazhari2024}         & \textcolor{symgreen}{\checkmark} & \textsc{\small S/B}  & \textcolor{symgreen}{\checkmark} & \textcolor{symgrey}{--}          & \textcolor{symgrey}{--}          & \textcolor{symgrey}{--}          \\
  \citet{JIANG2024101016}     & \textcolor{symgreen}{\checkmark} & \textcolor{symgrey}{--} & \textcolor{symgreen}{\checkmark} & \textcolor{symamber}{$\circ$} & \textcolor{symgrey}{--}          & \textcolor{symgrey}{--}          \\
  \citet{Nabot2025}           & \textcolor{symgreen}{\checkmark} & \textsc{\small Sc}   & \textcolor{symgreen}{\checkmark} & \textcolor{symgrey}{--}          & \textcolor{symgreen}{\checkmark} & \textcolor{symgreen}{\checkmark} \\
  \citet{Darvish2025}         & \textcolor{symgreen}{\checkmark} & \textsc{\small B}    & \textcolor{symgreen}{\checkmark} & \textcolor{symamber}{$\circ$}    & \textcolor{symamber}{$\circ$}    & \textcolor{symgreen}{\checkmark} \\
  \citet{Gholami2025}         & \textcolor{symgreen}{\checkmark} & \textsc{\small Sc}   & \textcolor{symgreen}{\checkmark} & \textcolor{symamber}{$\circ$}    & \textcolor{symgrey}{--}          & \textcolor{symgreen}{\checkmark} \\
  \citet{Kehinde2025}         & \textcolor{symgreen}{\checkmark} & \textsc{\small B}    & \textcolor{symgreen}{\checkmark} & \textcolor{symgreen}{\checkmark} & \textcolor{symgrey}{--}          & \textcolor{symgrey}{--}          \\
  \citet{Ajibade2025}         & \textcolor{symgreen}{\checkmark} & \textsc{\small Sc}   & \textcolor{symgreen}{\checkmark} & \textcolor{symamber}{$\circ$}    & \textcolor{symgrey}{--}          & \textcolor{symgrey}{--}          \\
  \citet{Huddar2025}          & \textcolor{symgreen}{\checkmark} & \textsc{\small B}    & \textcolor{symgreen}{\checkmark} & \textcolor{symamber}{$\circ$}    & \textcolor{symgrey}{--}          & \textcolor{symgrey}{--}          \\
  \citet{Akbari2026}          & \textcolor{symgreen}{\checkmark} & \textsc{\small B}    & \textcolor{symgreen}{\checkmark} & \textcolor{symamber}{$\circ$}    & \textcolor{symgreen}{\checkmark} & \textcolor{symgreen}{\checkmark} \\
  \citet{ebnerasoul2023}      & \textcolor{symgreen}{\checkmark} & \textsc{\small S}    & \textcolor{symgrey}{--}          & \textcolor{symamber}{$\circ$}    & \textcolor{symgrey}{--}          & \textcolor{symgrey}{--}          \\
  \citet{noroozi2016probabilistic}& \textcolor{symgreen}{\checkmark} & \textsc{\small B}  & \textcolor{symgrey}{--}          & \textcolor{symamber}{$\circ$}    & \textcolor{symgrey}{--}          & \textcolor{symgrey}{--}          \\
  \citet{JIANG2023104817}     & \textcolor{symgrey}{--}          & \textsc{\small B}    & \textcolor{symgreen}{\checkmark} & \textcolor{symgrey}{--}          & \textcolor{symgrey}{--}          & \textcolor{symgrey}{--}          \\
  \cmidrule{1-7}
  \rowcolor{pwbg}
    \color{headfg}\textbf{Present work}
    & \cellcolor{pwbg}\color{headfg}\checkmark
    & \color{headfg}\textbf{B}
    & \cellcolor{pwbg}\color{headfg}\checkmark
    & \color{headfg}\checkmark
    & \cellcolor{pwbg}\color{headfg}\checkmark
    & \cellcolor{pwbg}\color{headfg}\checkmark \\
  \bottomrule
\end{tabular}%
}
\end{table}